\documentclass{resume} % Use the custom resume.cls style
\usepackage[dvipsnames]{xcolor}
\usepackage{hyperref}
\DeclareUnicodeCharacter{2500}{---}

\usepackage[left=0.75in,top=0.6in,right=0.75in,bottom=0.1in]{geometry} % Document margins
\newcommand{\tab}[1]{\hspace{.2667\textwidth}\rlap{#1}}
\newcommand{\itab}[1]{\hspace{0em}\rlap{#1}}
\newcommand{\myname}{\textcolor{CarnegieMellonRed}{\textbf{Z. Wang}}}
\newcommand{\pubentry}[6]{\noindent[#1] #2, #3, \textit{#4}, #5 \href{https://arxiv.org/abs/#6}{[#6]}\par}
\name{Zihao(Matthew) Wang} % Your name
\address{Homepage: \href{https://zihaowang-kiaa.github.io/personal-homepage/}{\texttt{website}}}
\address{
\href{https://github.com/zihaowang-kiaa}{Github} \\
\href{https://orcid.org/my-orcid?orcid=0009-0007-4590-6550}{Orcid} \\
\href{ https://ui.adsabs.harvard.edu/public-libraries/KrWPK55_RIi-0et1K3-kDQ}{NASA ADS}
}

\address{
Email: \texttt{zihaowang25@stu.pku.edu.cn} \\
Contact: (+86) 13309023239
}
% Your phone number and email

\definecolor{CarnegieMellonRed}{RGB}{196,18,48}

\renewenvironment{rSection}[1]{
\sectionskip
\textcolor{CarnegieMellonRed}{\MakeUppercase{#1}}
\sectionlineskip
\hrule
\begin{list}{}{
\setlength{\leftmargin}{1.5em}
}
\item[]
}{
\end{list}
}



\begin{document}
\begin{rSection}{Education and Experience}
{\bf Peking University} \hfill {\em Sept 2025 - June 2030 (expected)}
\\ PhD student in Astronomy \hfill
\\  Advisor: Prof. Kohei Inayoshi \hfill
% \\ \textit{Related courses: A, B, C. }

{\bf Nanjing University} \hfill {\em Sept 2021 - June 2025}
\\ B.S. in Astronomy (Advanced Class)
\\ \qquad GPA: 4.54/5 (Rank: 1/ 42)\qquad  \hfill
\\ \qquad Thesis: \textit{Estimation and physical origin of high-redshift galaxy properties in the JWST era}
\\  Advisor: Prof. Tao Wang \hfill

{\bf Massachusetts Institute of Technology} \hfill {\em June 2024 - Decem 2024}
\\ Visiting research student in Physics
\\  Advisor: Dr. Xuejian Shen \& Prof. Mark Vogelsberger\hfill
\end{rSection}

\begin{rSection}{Research Interests}
{\bf High-redshift AGN \& Little Red Dots:} Theoretical modeling and predictions of the origin and evolution of Little Red Dots; Large-scale environments of high-$z$ AGN; High-$z$ AGN population in cosmological simulations \hfill

{\bf High-redshift galaxies \& star formation:} Multi-scale star formation at high-$z$ from cloud and cluster to galactic scales; Diverse morphology of high-$z$ galaxies; Metal enrichment \hfill

{\bf Machine learning in astrophysics:} Use of machine learning techniques in galaxy parameter estimation, classification, and target hunting. \hfill


\end{rSection}


% \begin{rSection}{Research Experience}
% \begin{rSubsection}{Project name 1 \cite{paper1}}{Jan 2024 - Present}{Supervisors: \href{https://example.com/}{Dr. A}, \href{https://example.com/}{Prof. B}, and \href{https://example.com/}{Prof. C}}{CMU}
% \item Proposed ...
% \item Found ...
% \item Completed ...
% \item Modified ...
% \item Proposed ...
% \item Finished ...
% \end{rSubsection}


% \begin{rSubsection}{Project name 2 \cite{paper2}}{Dec 2023 - Feb 2024}{Supervisors: \href{https://example.weebly.com/}{Dr. A} and \href{https://www.ece.cmu.edu/directory/bios/example}{Prof. B}}{CMU}
% \item Ran experiments for ...
% \item Organized ...
% \item Explored ...
% \item Proposed ...
% \item Finished ...
% \end{rSubsection}
% \end{rSection}


\begin{rSection}{Publications} \itemsep -2pt
% {\bf Selected publications.} \textcolor{CarnegieMellonRed}{\textbf{Z. Wang}} indicates my authorship.\par\smallskip
{\bf Publications (as of 2026/8): 4 first-author (2 peer-reviewed) and 3 contributing-author (2 peer-reviewed) papers;
89 total citations.}

{\bf First-author papers}\par
\pubentry{1}{\myname, X. Shen et al.}{The THESAN-ZOOM project: clumpiness of high-redshift galaxies and its connection to bursty star formation}{arXiv}{2026}{2608.19308}
\pubentry{2}{\myname, F. Jiang et al.}{Halo assembly bias in the early Universe: a clustering probe of the origin of the Little Red Dots}{arXiv}{2026}{2603.15736}
\pubentry{3}{\myname, T. Wang et al.}{SHAPE: I. A SOM-SED hybrid approach for efficient galaxy parameter estimation leveraging JWST}{A\&A}{2026}{2510.00187}
\pubentry{4}{\myname, X. Shen et al.}{The THESAN-ZOOM project: star formation efficiency from giant molecular clouds to galactic scale in high-redshift starbursts}{MNRAS}{2025}{2505.05554}

{\bf Contributing-author papers}\par
\pubentry{5}{W. McClymont, V. Belokurov et al., including \myname}{Globular cluster abundance patterns inherited from giant molecular clouds}{arXiv}{2026}{2607.05509}
\pubentry{6}{W. McClymont, S. Tacchella et al., including \myname}{The THESAN-ZOOM project: mystery N/O more --- uncovering the origin of peculiar chemical abundances and a not-so-fundamental metallicity relation at $3<z<12$}{MNRAS}{2026}{2507.08787}
\pubentry{7}{X. Shen, R. Kannan et al., including \myname}{The THESAN-ZOOM project: star formation efficiencies in high-redshift galaxies}{MNRAS}{2026}{2503.01949}

\end{rSection}

% \newpage

\begin{rSection}{Talks and Presentations} \itemsep -2pt
{[1] Conference Talk, Guoshoujing annual meeting, Haikou, China} \\
\textit{Where Little Red Dots live tells us what they are} \hfill {\em 2026}


{[2] Conference Talk, CSST annual meeting, Kunming, China} \\
\textit{Introducing SHAPE: A SOM-SED Hybrid Approach for Efficient galaxy Parameter Estimation}   \hfill {\em 2025}

{[3] Conference Poster, Sino-German conference on Galaxy Formation, Chengdu, China} \\
\textit{Elevated star formation efficiency at high-$z$: from a giant molecular cloud perspective} \hfill {\em 2025}

\end{rSection}

\begin{rSection}{Honors and awards} \itemsep -2pt
{Outstanding Undergraduate Thesis of Nanjing University (the special prize)}\hfill {\em 2025} \\
{Zheng-Gang Oversea Scholarship of Nanjing University}\hfill {\em 2024} \\
{Elite Program Scholarship (the special prize)}\hfill {\em 2024} \\
{People’s Scholarship (the first prize)} \hfill {\em 2022, 2023, 2024} \\
{National Astronomical Observatories Scholarship}\hfill {\em 2022}
\end{rSection}



\begin{rSection}{Skills/Hobbies (* denotes extensively)} \itemsep -2pt
\begin{tabular}{ @{} >{\bfseries}l @{\hspace{6ex}} l }
Programming Languages &  Python*, C/C++, HTML \\
Machine Learning Tools & SOM*, CNN, Random-Forest \\
Language & Mandarin (Native), English (Proficient) \\
Simulations &  \textit{Thesan-Zoom}*, \textit{LUMINA}*, \textit{ShinUchuu}, \textit{IllustrisTNG}, \textit{FIRE} \\
Hobbies & cinema and photographing \\
\end{tabular}
\end{rSection}


\begin{rSection}{Professional references} \itemsep -2pt
{\bf Prof. Kohei Inayoshi } \\
Kavli Institute for Astronomy and Astrophysics, PKU \hfill {Email: \texttt{inayoshi0328@gmail.com}}

{\bf Prof. Fangzhou Jiang } \\
Kavli Institute for Astronomy and Astrophysics, PKU \hfill {Email: \texttt{fangzhou.jiang@pku.edu.cn}}

{\bf Dr. Xuejian Shen} \\
Kavli Institute for Astrophysics and Space Research, MIT \hfill {Email: \texttt{xuejian@mit.edu}}


{\bf Prof. Mark Vogelsberger} \\
Kavli Institute for Astrophysics and Space Research, MIT \hfill {Email: \texttt{mvogelsb@mit.edu}}

{\bf Prof. Tao Wang} \\
School of Astronomy and Space Science, NJU
\hfill {Email: \texttt{taowang@nju.edu.cn}}


\end{rSection}





\end{document}
